\documentclass[a4paper,11pt]{ctexart}
\usepackage{amsmath, amssymb, amsfonts}
\usepackage{geometry}
\geometry{left=1.5cm, right=1.5cm, top=2.0cm, bottom=2.0cm, headsep=0.3cm, footskip=1cm}
\usepackage{listings}
\usepackage{xcolor}
\usepackage{tikz}
\usepackage{pgfplots}
\usepackage{hyperref}
\usepackage{fancyhdr}
\usepackage{tcolorbox}
\tcbuselibrary{breakable, skins}
\usepackage{array}
\usepackage{booktabs}
\usepackage[shortlabels]{enumitem}
\usepackage{needspace}
\usepackage{multirow}

\AtBeginDocument{
  \hypersetup{colorlinks=true, linkcolor=blue, filecolor=magenta, urlcolor=cyan}
}

\setlist{nosep, itemsep=0.8ex, parsep=0.1em}
\setlist[itemize]{leftmargin=1.8em}
\setlist[enumerate]{leftmargin=1.8em}

% ===== 颜色定义 =====
\definecolor{codegreen}{rgb}{0,0.6,0}
\definecolor{codegray}{rgb}{0.5,0.5,0.5}
\definecolor{codepurple}{rgb}{0.58,0,0.82}
\definecolor{codebg}{rgb}{0.98,0.98,0.98}
\definecolor{tipsblue}{rgb}{0.85,0.93,1.0}
\definecolor{highlightyellow}{rgb}{1.0,0.97,0.8}

% ===== 代码块样式 =====
\lstdefinestyle{mystyle}{
    backgroundcolor=\color{codebg},
    commentstyle=\color{codegreen},
    keywordstyle=\color{magenta}\bfseries,
    numberstyle=\tiny\color{codegray},
    stringstyle=\color{codepurple},
    basicstyle=\ttfamily\small,
    breakatwhitespace=false,
    breaklines=true,
    captionpos=b,
    keepspaces=true,
    numbers=left,
    numbersep=6pt,
    showspaces=false,
    showstringspaces=false,
    showtabs=false,
    tabsize=2,
    language=Python,
    frame=single,
    framerule=0.5pt,
    rulecolor=\color{gray!40}
}
\lstset{style=mystyle}

% ===== 彩色框 =====
\newtcolorbox{problembox}[1][]{
    colback=white,
    colframe=blue!60!black,
    title=\textbf{#1},
    fonttitle=\bfseries\small,
    boxrule=1pt,
    arc=3pt, outer arc=3pt,
    coltitle=black,
    before skip=10pt, after skip=8pt
}

\newtcolorbox{solutionbox}[1][]{
    enhanced, breakable,
    colback=green!3!white,
    colframe=green!50!black,
    title=\textbf{#1},
    fonttitle=\bfseries\small,
    boxrule=1pt,
    arc=3pt, outer arc=3pt,
    coltitle=black,
    before skip=8pt, after skip=10pt
}

\newtcolorbox{metaphorbox}[1][]{
    colback=orange!5!white,
    colframe=orange!70!black,
    title=\textbf{#1},
    fonttitle=\bfseries\small,
    boxrule=1pt,
    arc=3pt, outer arc=3pt,
    coltitle=black,
    before skip=8pt, after skip=8pt
}

\newtcolorbox{beginnerbox}[1][]{
    colback=tipsblue,
    colframe=blue!50!black,
    title=\textbf{#1},
    fonttitle=\bfseries\small,
    boxrule=1pt,
    arc=3pt, outer arc=3pt,
    coltitle=black,
    before skip=8pt, after skip=8pt
}

\newtcolorbox{appbox}[1][]{
    colback=green!3!white,
    colframe=green!60!black,
    title=\textbf{#1},
    fonttitle=\bfseries\small,
    boxrule=1pt,
    arc=3pt, outer arc=3pt,
    coltitle=black,
    before skip=8pt, after skip=10pt
}

\newtcolorbox{keybox}[1][]{
    colback=highlightyellow,
    colframe=orange!80!black,
    title=\textbf{#1},
    fonttitle=\bfseries\small,
    boxrule=1pt,
    arc=3pt, outer arc=3pt,
    coltitle=black,
    before skip=8pt, after skip=8pt
}

% ===== 页眉页脚 =====
\pagestyle{fancy}
\fancyhf{}
\fancyhead[R]{深度学习-第5章}
\fancyhead[L]{学习的技巧 · 练习题}
\fancyfoot[C]{\thepage}
\addtolength{\headheight}{2pt}

\parskip=2pt plus 1pt minus 0.5pt
\linespread{1.35}
\raggedbottom
\widowpenalty=10000
\clubpenalty=10000

\title{\textbf{深度学习\\第5章\quad 学习的技巧}\\[0.5em]
{\large\color{gray}——梯度、初始化、正则化：让深层网络真正"学得起来"的三把钥匙}}
\date{\today}

\begin{document}

\maketitle

% ===== 章节导读 =====
\begin{beginnerbox}[本章题目导览：你将挑战哪些核心问题？]
    搭一个深层神经网络很容易，但让它真正\textbf{稳定、高效地训练}却困难得多。梯度会消失或爆炸，优化器会原地震荡，权重初始化不当会让网络一开始就"跑偏"，过拟合会让高分训练集变成低分测试集。本章 6 道题逐一击破这些拦路虎：

    \begin{itemize}
        \item \textbf{Q5-01（辨析填空）}：梯度消失 vs 梯度爆炸——深层网络为什么容易出现这两种问题？各自的症状、根源和解决思路是什么？
        \item \textbf{Q5-02（计算题）}：优化器更新规则手动推导——给定梯度历史，分别用 SGD、Momentum、AdaGrad 手算一步参数更新，把公式变成具体数字。
        \item \textbf{Q5-03（对比分析）}：六种优化器横向对比——从 SGD 到 Adam，各自解决了什么问题、引入了什么超参数、适合什么场景？
        \item \textbf{Q5-04（概念理解）}：六种参数初始化方法对比——常数、秩、正态、均匀、Xavier、He 初始化，各自的公式和适用激活函数是什么？
        \item \textbf{Q5-05（代码阅读）}：Batch Normalization 代码追踪——BN 层的计算过程、训练/推理行为差异，以及在网络中的位置选择。
        \item \textbf{Q5-06（综合分析）}：过拟合三剑客——权值衰减、Dropout、BatchNorm 各自如何抑制过拟合？它们能叠加使用吗？
    \end{itemize}
\end{beginnerbox}

\section{第 5 章\quad 学习的技巧}

前面几章我们已经能搭出一个神经网络并跑通训练循环了。但"能跑"和"能学好"之间，还差着一大截：

\begin{itemize}
    \item 网络加深到10层，梯度传到底层已经变成接近0的小数——参数根本没在更新；
    \item SGD 在峡谷形损失曲面上左右震荡，迟迟无法到达谷底；
    \item 初始权重太大或太小，激活值要么全饱和、要么全为零——无论怎么训练都没有起色；
    \item 训练集准确率99\%，验证集准确率65\%——模型把噪声也背下来了。
\end{itemize}

这些不是玄学问题，都有对应的\textbf{技巧}来解决。本章就是这本"工具手册"。

\begin{metaphorbox}[先建立一个总感觉：为什么深度学习需要"技巧"？]
    \textbf{把训练深层神经网络比作在山脉中寻找最低谷。}

    \begin{itemize}
        \item \textbf{梯度消失/爆炸} = 山太高太深，向导（梯度）的声音传到山底时已经完全失真——要么细如蚊鸣（消失），要么震耳欲聋（爆炸），两种情况下你都找不到方向；
        \item \textbf{优化器} = 你的行走策略：原地踏步（SGD 震荡）、借助惯性滑行（Momentum）、根据地形自动调节步伐（AdaGrad/Adam）；
        \item \textbf{参数初始化} = 出发点的选择：从悬崖边出发（权重过大）或者从坑底出发（权重过小），都很难找到最优路线；
        \item \textbf{正则化} = 防止你把地图上的噪声当成真实地形记住——记住规律，而不是死记硬背。
    \end{itemize}
\end{metaphorbox}


%% ============================================================
\Needspace{5\baselineskip}
\begin{problembox}[Q5-01\quad 梯度消失 vs 梯度爆炸：深层网络的两大顽疾]
    \textbf{题目描述：}\\
    文档第 5.1 节指出，深度神经网络在训练过程中面临梯度消失（Vanishing Gradient）和梯度爆炸（Exploding Gradient）两大问题。两者都与反向传播中的链式法则密切相关。

    \vspace{0.3cm}
    \textbf{问题：}
    \begin{enumerate}[(1)]
        \item 填写下表，对比两种问题的核心特征：

        \begin{center}
        \begin{tabular}{p{2.5cm} p{4.5cm} p{4.5cm}}
            \toprule
            & \textbf{梯度消失} & \textbf{梯度爆炸} \\
            \midrule
            根本原因 & 链式法则连乘时，每层梯度值\underline{\hspace{1.5cm}} & 链式法则连乘时，每层梯度值\underline{\hspace{1.5cm}} \\[1.5em]
            典型症状 & 底层参数\underline{\hspace{2.0cm}}，Loss 不再下降 & 参数更新幅度极大，Loss\underline{\hspace{1.5cm}}或变成 NaN \\[1.5em]
            易发激活函数 & \underline{\hspace{3.0cm}} & 通常与激活函数\underline{\hspace{1.5cm}}，与权重初始化更相关 \\[1.5em]
            常见解决方案 & 换用\underline{\hspace{1.5cm}}激活函数；使用\underline{\hspace{1.5cm}}初始化 & 梯度裁剪（Gradient Clipping）；使用\underline{\hspace{1.5cm}}初始化 \\
            \bottomrule
        \end{tabular}
        \end{center}

        \item \textbf{数学推导}：设一个5层网络，每层的梯度（局部导数）分别为 $g_1, g_2, g_3, g_4, g_5$。根据链式法则，第1层参数的梯度为：
        \[
        \frac{\partial L}{\partial w_1} = g_5 \cdot g_4 \cdot g_3 \cdot g_2 \cdot g_1
        \]
        若每层使用 Sigmoid 激活函数，导数最大值约为 $0.25$，计算5层连乘后第1层的梯度量级（假设每层导数恰好取最大值 $0.25$）。

        \item \textbf{ReLU 为何能缓解梯度消失}：ReLU 在正值区域的导数为常数1。解释为什么这能让梯度"畅通"地传递到底层，以及 ReLU 仍然可能引发梯度问题的场景（死亡 ReLU）。

        \item \textbf{梯度裁剪}是缓解梯度爆炸最直接的方法，其核心思想是：当梯度的范数超过阈值 $\theta$ 时，对梯度进行等比例缩放：
        \[
        \mathbf{g} \leftarrow \frac{\theta}{\|\mathbf{g}\|} \cdot \mathbf{g}, \quad \text{当 } \|\mathbf{g}\| > \theta
        \]
        若当前梯度向量 $\mathbf{g} = [6, 8]$，裁剪阈值 $\theta = 5$，计算裁剪后的梯度。
    \end{enumerate}

    \vspace{0.2cm}
    {\color{gray}\small\textit{来源溯源：[文档第 5 章 5.1 深度神经网络及其问题]}}
\end{problembox}

\begin{beginnerbox}[小白必读：反向传播为什么会让梯度"越传越小"？]
    \begin{itemize}
        \item \textbf{链式法则的连乘效应}：反向传播本质上是从输出层向输入层逐层相乘梯度。如果每层的梯度值都是一个小于1的小数（比如 Sigmoid 导数最大只有0.25），乘了10层之后，梯度就变成了 $0.25^{10} \approx 10^{-6}$，几乎为零，底层参数收到的"纠错信号"消失殆尽。
        \item \textbf{梯度爆炸的对称情形}：如果每层梯度大于1（比如权重初始化过大，导致每层放大梯度），连乘后梯度急剧增大，参数更新幅度失控，Loss 剧烈震荡甚至变成 \texttt{NaN}。
        \item \textbf{深度越深，问题越严重}：层数翻倍，连乘次数翻倍，梯度消失/爆炸的程度呈指数级恶化。这正是早期深度网络（10层以上）难以训练的核心原因。
    \end{itemize}
\end{beginnerbox}

\begin{solutionbox}[参考答案与详细解析]
    \textbf{(1) 对比表填写：}

    \begin{center}
    \begin{tabular}{p{2.5cm} p{4.5cm} p{4.5cm}}
        \toprule
        & \textbf{梯度消失} & \textbf{梯度爆炸} \\
        \midrule
        根本原因 & 连乘时每层梯度值小于1，连乘后趋近于0 & 连乘时每层梯度值大于1，连乘后趋向无穷 \\[1em]
        典型症状 & 底层参数几乎不更新，Loss 不再下降 & Loss 剧烈震荡或变成 NaN \\[1em]
        易发激活函数 & Sigmoid、Tanh（导数最大值$<$1） & 关系较小，主要与权重初始化相关 \\[1em]
        常见解决方案 & 换用 ReLU；使用 Xavier/He 初始化；BatchNorm & 梯度裁剪；使用合理初始化 \\
        \bottomrule
    \end{tabular}
    \end{center}

    \textbf{(2) 梯度量级计算：}

    假设每层 Sigmoid 导数取最大值 $0.25$，5层连乘：
    \[
    \frac{\partial L}{\partial w_1} = 0.25^5 = \frac{1}{4^5} = \frac{1}{1024} \approx \boxed{9.77 \times 10^{-4}}
    \]
    不到千分之一。若网络有10层，则约为 $0.25^{10} \approx 9.5 \times 10^{-7}$，几乎无法更新。

    \textbf{(3) ReLU 缓解梯度消失的原因：}

    ReLU 在正值区域导数恒为1，连乘后梯度不会因层数增加而衰减，信号可以畅通传递。

    \textbf{死亡 ReLU（Dying ReLU）}：当某个神经元的输入永远为负时（如参数更新后该神经元输出始终$\leq 0$），ReLU 输出恒为0，导数恒为0，这个神经元永远不再接收梯度更新——它"死亡"了。解决方案：使用 Leaky ReLU（负值区域有小斜率 $\alpha x$，而非直接为0）或调小学习率避免参数大幅跑偏。

    \textbf{(4) 梯度裁剪计算：}

    当前梯度 $\mathbf{g} = [6, 8]$，范数 $\|\mathbf{g}\| = \sqrt{6^2 + 8^2} = \sqrt{100} = 10$。

    由于 $10 > \theta = 5$，触发裁剪：
    \[
    \mathbf{g}_{\text{clip}} = \frac{5}{10} \cdot [6, 8] = \boxed{[3, 4]}
    \]
    裁剪后范数恰好等于阈值5，方向不变（仍为 $[6,8]$ 的方向），幅度被等比缩小。
\end{solutionbox}


%% ============================================================
\Needspace{5\baselineskip}
\begin{problembox}[Q5-02\quad 优化器更新规则手动推导]
    \textbf{题目描述：}\\
    文档第 5.2 节依次介绍了 SGD、Momentum、AdaGrad 等优化器的更新规则。理解公式最好的方式是亲手算一遍。

    \vspace{0.3cm}
    \textbf{场景设定：}某参数 $w$ 当前值为 $w_0 = 1.0$，学习率 $\eta = 0.1$。在连续两步中，梯度分别为 $g_1 = 0.5$（第1步）和 $g_2 = 0.8$（第2步）。

    \vspace{0.3cm}
    \textbf{问题：}
    \begin{enumerate}[(1)]
        \item \textbf{标准 SGD} 更新规则为 $w \leftarrow w - \eta \cdot g$。计算两步后 $w$ 的值。

        \item \textbf{Momentum} 更新规则为：
        \[
        v \leftarrow \beta v - \eta \cdot g, \quad w \leftarrow w + v
        \]
        取 $\beta = 0.9$，初始动量 $v_0 = 0$。分别计算第1步、第2步后的 $v$ 和 $w$（保留三位小数）。

        \item \textbf{AdaGrad} 更新规则为：
        \[
        h \leftarrow h + g^2, \quad w \leftarrow w - \frac{\eta}{\sqrt{h} + \epsilon} \cdot g
        \]
        取 $\epsilon = 10^{-8}$（防止除零，计算时可忽略），初始 $h_0 = 0$。计算第1步、第2步后的 $h$ 和 $w$（保留三位小数）。

        \item 对比三种方法在两步之后的 $w$ 值，并分析：为什么在第2步（梯度更大）时，AdaGrad 的更新幅度反而比 SGD \textbf{更小}？这反映了 AdaGrad 的什么特性？
    \end{enumerate}

    \vspace{0.2cm}
    {\color{gray}\small\textit{来源溯源：[文档第 5 章 5.2 更新参数方法的优化]}}
\end{problembox}

\begin{beginnerbox}[小白必读：为什么 SGD 的缺点那么明显却还在用？]
    \begin{itemize}
        \item \textbf{SGD 的核心缺陷}：对所有参数使用同一个学习率，在损失曲面各方向曲率差异很大时（"峡谷"形曲面），在陡峭方向步子太大导致震荡，在平坦方向步子太小导致收敛极慢。
        \item \textbf{为什么还在用}：在大型模型（如 ResNet、GPT）训练中，经过仔细调学习率和配合学习率衰减后，SGD+Momentum 有时比 Adam 收敛到更好的泛化性能——这在计算机视觉领域尤为明显。Adam 虽然收敛快，但有时训练集Loss 下降很快而测试集表现不如 SGD+Momentum。
        \item \textbf{实用建议}：入门用 Adam（调参简单，效果不错）；需要极致性能时，尝试 SGD+Momentum+学习率衰减。
    \end{itemize}
\end{beginnerbox}

\begin{solutionbox}[参考答案与详细解析]
    \textbf{(1) 标准 SGD：}

    \textbf{第1步}（$g_1 = 0.5$）：
    \[
    w_1 = 1.0 - 0.1 \times 0.5 = 1.0 - 0.05 = 0.95
    \]
    \textbf{第2步}（$g_2 = 0.8$）：
    \[
    w_2 = 0.95 - 0.1 \times 0.8 = 0.95 - 0.08 = \boxed{0.87}
    \]

    \textbf{(2) Momentum（$\beta=0.9$，$v_0=0$）：}

    \textbf{第1步}：
    \[
    v_1 = 0.9 \times 0 - 0.1 \times 0.5 = -0.05
    \]
    \[
    w_1 = 1.0 + (-0.05) = 0.95
    \]
    \textbf{第2步}：
    \[
    v_2 = 0.9 \times (-0.05) - 0.1 \times 0.8 = -0.045 - 0.08 = -0.125
    \]
    \[
    w_2 = 0.95 + (-0.125) = \boxed{0.825}
    \]
    Momentum 将历史梯度的"惯性"（$v_1=-0.05$）叠加到当前更新上，步子比 SGD 更大，收敛更快。

    \textbf{(3) AdaGrad（$\epsilon$ 忽略，$h_0=0$）：}

    \textbf{第1步}（$g_1=0.5$）：
    \[
    h_1 = 0 + 0.5^2 = 0.25
    \]
    \[
    w_1 = 1.0 - \frac{0.1}{\sqrt{0.25}} \times 0.5 = 1.0 - \frac{0.1}{0.5} \times 0.5 = 1.0 - 0.1 = 0.9
    \]
    \textbf{第2步}（$g_2=0.8$）：
    \[
    h_2 = 0.25 + 0.8^2 = 0.25 + 0.64 = 0.89
    \]
    \[
    w_2 = 0.9 - \frac{0.1}{\sqrt{0.89}} \times 0.8 = 0.9 - \frac{0.1}{0.9434} \times 0.8 \approx 0.9 - 0.0848 = \boxed{0.815}
    \]

    \textbf{(4) 对比分析：}

    \begin{center}
    \begin{tabular}{lccc}
        \toprule
        优化器 & $w_0$ & $w_1$ & $w_2$ \\
        \midrule
        SGD & 1.000 & 0.950 & 0.870 \\
        Momentum & 1.000 & 0.950 & 0.825 \\
        AdaGrad & 1.000 & 0.900 & 0.815 \\
        \bottomrule
    \end{tabular}
    \end{center}

    在第2步，SGD 更新幅度为 $0.08$，AdaGrad 约为 $0.085$（此例中 $h$ 的累积使有效学习率约为 $0.106$，但分母中的 $\sqrt{h}$ 已大于1）。随着训练步数增加，$h$ 持续累积增大，分母 $\sqrt{h}$ 越来越大，有效学习率 $\eta/\sqrt{h}$ 单调递减——\textbf{历史梯度越大的参数，后续学习率越小}。这正是 AdaGrad 的"自适应学习率"特性：频繁更新的参数自动减小步长，稀疏梯度的参数维持大步长，特别适合稀疏特征数据（如词频）。

    \begin{keybox}[AdaGrad 的致命弱点]
        $h$ 是\textbf{单调递增}的——它只累加，从不遗忘。训练到后期，$h$ 极大，有效学习率趋近于0，参数几乎无法更新。RMSProp 用\textbf{指数移动平均}代替全量累加，解决了这个问题。
    \end{keybox}
\end{solutionbox}


%% ============================================================
\Needspace{5\baselineskip}
\begin{problembox}[Q5-03\quad 六种优化器横向对比]
    \textbf{题目描述：}\\
    文档第 5.2 节完整介绍了从 SGD 到 Adam 的演进历程。每一种优化器的诞生，都是为了解决上一种的某个缺点。

    \vspace{0.3cm}
    \textbf{问题：}
    \begin{enumerate}[(1)]
        \item 填写下表，梳理六种优化器的核心信息：

        \begin{center}
        \begin{tabular}{p{2.2cm} p{4.0cm} p{2.5cm} p{3.5cm}}
            \toprule
            优化器 & 核心改进点 & 关键超参数 & 适用场景 \\
            \midrule
            SGD & 基础梯度下降，无改进 & 学习率 $\eta$ & \underline{\hspace{3cm}} \\[1.2em]
            Momentum & 引入\underline{\hspace{2.0cm}}，赋予惯性 & $\eta$，动量系数 $\beta$ & \underline{\hspace{3cm}} \\[1.2em]
            学习率衰减 & 训练后期\underline{\hspace{2.0cm}}，避免震荡 & 衰减策略、衰减率 & \underline{\hspace{3cm}} \\[1.2em]
            AdaGrad & 自适应学习率，历史梯度平方\underline{\hspace{1.5cm}} & $\eta$，$\epsilon$ & \underline{\hspace{3cm}} \\[1.2em]
            RMSProp & 改用\underline{\hspace{2.0cm}}，解决 AdaGrad 学习率归零 & $\eta$，衰减率 $\rho$ & \underline{\hspace{3cm}} \\[1.2em]
            Adam & 融合\underline{\hspace{1cm}}和\underline{\hspace{1cm}}，加偏差修正 & $\eta$，$\beta_1$，$\beta_2$，$\epsilon$ & \underline{\hspace{3cm}} \\
            \bottomrule
        \end{tabular}
        \end{center}

        \item \textbf{学习率衰减策略}：列举两种常见的学习率衰减方式，写出其调度规则，并说明各自的优缺点。

        \item \textbf{Adam 的两个动量}：Adam 同时维护一阶动量（梯度的指数移动平均 $m_t$）和二阶动量（梯度平方的指数移动平均 $v_t$），其更新公式为：
        \[
        m_t = \beta_1 m_{t-1} + (1-\beta_1) g_t, \quad v_t = \beta_2 v_{t-1} + (1-\beta_2) g_t^2
        \]
        \[
        \hat{m}_t = \frac{m_t}{1-\beta_1^t}, \quad \hat{v}_t = \frac{v_t}{1-\beta_2^t}, \quad w \leftarrow w - \frac{\eta}{\sqrt{\hat{v}_t}+\epsilon}\hat{m}_t
        \]
        解释一阶动量和二阶动量各自的"功能"，以及为什么需要偏差修正（$\hat{m}_t$、$\hat{v}_t$）。

        \item 在深度学习实践中，如何选择优化器？给出一个决策流程建议（考虑：任务类型、数据集大小、调参成本等因素）。
    \end{enumerate}

    \vspace{0.2cm}
    {\color{gray}\small\textit{来源溯源：[文档第 5 章 5.2 更新参数方法的优化]}}
\end{problembox}

\begin{metaphorbox}[生活比喻：六种优化器像六种下山方式]
    你站在山顶，要找到最低点，脚下只能感受到坡度（梯度）：
    \begin{itemize}
        \item \textbf{SGD}：每步严格朝当前最陡方向迈固定步伐，遇到峡谷就来回震荡；
        \item \textbf{Momentum}：像滚石，有惯性，能冲出浅坑，但可能冲过头；
        \item \textbf{学习率衰减}：开始步子大，接近谷底时主动放慢——像跑步冲线前收速；
        \item \textbf{AdaGrad}：走过陡路（大梯度）后主动缩小步伐，但步子只会越来越小；
        \item \textbf{RMSProp}：和 AdaGrad 一样聪明，但只记最近的路况，步子不会无限缩小；
        \item \textbf{Adam}：既有滚石的惯性（$m_t$），又能根据近期地形自动调步幅（$v_t$），综合最优，但在某些地形下可能错过"更好的最低点"。
    \end{itemize}
\end{metaphorbox}

\begin{solutionbox}[参考答案与详细解析]
    \textbf{(1) 六种优化器对比表：}

    \begin{center}
    \begin{tabular}{p{2.2cm} p{4.0cm} p{2.5cm} p{3.0cm}}
        \toprule
        优化器 & 核心改进点 & 关键超参数 & 适用场景 \\
        \midrule
        SGD & 基础梯度下降 & $\eta$ & 凸问题/配合精细调参 \\[0.8em]
        Momentum & 引入历史梯度惯性 & $\eta$，$\beta$ & 损失曲面有"峡谷"时 \\[0.8em]
        学习率衰减 & 训练后期自动降低学习率 & 衰减策略/衰减率 & 所有场景的辅助策略 \\[0.8em]
        AdaGrad & 历史梯度平方累加 & $\eta$，$\epsilon$ & 稀疏特征/NLP词频 \\[0.8em]
        RMSProp & 指数移动平均代替全量累加 & $\eta$，$\rho$ & 非平稳目标/RNN \\[0.8em]
        Adam & 融合Momentum和RMSProp & $\eta$，$\beta_1$，$\beta_2$，$\epsilon$ & 大多数深度学习任务 \\
        \bottomrule
    \end{tabular}
    \end{center}

    \textbf{(2) 两种学习率衰减策略：}

    \begin{itemize}
        \item \textbf{步进衰减（Step Decay）}：每隔固定步数/epoch，将学习率乘以衰减因子 $\gamma$（如每10个epoch乘以0.1）。优点：简单直观，易于控制；缺点：学习率阶梯式突变，可能引起短暂震荡。
        \item \textbf{余弦退火（Cosine Annealing）}：学习率按余弦函数从初始值平滑降至接近0：$\eta_t = \eta_{\min} + \frac{1}{2}(\eta_{\max}-\eta_{\min})(1+\cos\frac{t\pi}{T})$。优点：过渡平滑，最终阶段学习率极小有利于精细收敛；缺点：超参数（$T$，$\eta_{\min}$）需要根据训练总步数预先设定。
    \end{itemize}

    \textbf{(3) Adam 两个动量的功能：}

    \begin{itemize}
        \item \textbf{一阶动量 $m_t$}（梯度的移动平均）：扮演 Momentum 的角色，记录梯度的方向趋势，赋予更新"惯性"，帮助跨越局部极小值；
        \item \textbf{二阶动量 $v_t$}（梯度平方的移动平均）：扮演 RMSProp 的角色，估计梯度的幅度，为不同参数自适应调节步长——历史梯度大的参数步长小，梯度小的步长大。
    \end{itemize}

    \textbf{偏差修正的必要性}：初始时 $m_0=v_0=0$，前几步的 $m_t$、$v_t$ 会被初始零值"拉低"，严重低估真实动量。修正项 $1/(1-\beta^t)$ 在初始阶段（$t$ 小）较大，对估计值进行放大，还原真实动量；随着训练进行（$t$ 增大），$\beta^t \to 0$，修正项趋近于1，影响自然消失。

    \textbf{(4) 优化器选择建议：}

    \begin{itemize}
        \item \textbf{默认首选 Adam}（$\beta_1=0.9$，$\beta_2=0.999$，$\eta=10^{-3}$）：适合绝大多数任务，调参成本低，收敛快；
        \item \textbf{计算机视觉分类任务}：可以尝试 SGD + Momentum + 余弦退火，有时泛化更好；
        \item \textbf{NLP/稀疏特征}：AdaGrad 或 Adam 均可；Transformer 训练通常用 AdamW（Adam + 权重衰减解耦）；
        \item \textbf{资源受限/快速验证}：直接用 Adam，先跑通再优化；
        \item \textbf{训练不稳定时}：降低学习率，加入学习率预热（Warmup），或切换为 SGD+Momentum。
    \end{itemize}
\end{solutionbox}


%% ============================================================
\Needspace{5\baselineskip}
\begin{problembox}[Q5-04\quad 六种参数初始化方法对比]
    \textbf{题目描述：}\\
    文档第 5.3 节系统介绍了六种参数初始化方法。初始化是训练的"第一步"，选错了，后面再好的优化器也难以弥补。

    \vspace{0.3cm}
    \textbf{问题：}
    \begin{enumerate}[(1)]
        \item 填写下表，梳理六种初始化方法的关键信息：

        \begin{center}
        \begin{tabular}{p{2.8cm} p{4.5cm} p{4.0cm}}
            \toprule
            初始化方法 & 初始化规则 & 适用场景/注意事项 \\
            \midrule
            常数初始化 & 所有权重初始化为常数 $c$ & \underline{\hspace{3.5cm}} \\[1.2em]
            秩初始化 & 使用\underline{\hspace{2.5cm}}矩阵分解确定初值 & \underline{\hspace{3.5cm}} \\[1.2em]
            正态分布初始化 & $W \sim \mathcal{N}(0,\ \sigma^2)$ & $\sigma$ 过大易\underline{\hspace{1.5cm}}，过小易\underline{\hspace{1.5cm}} \\[1.2em]
            均匀分布初始化 & $W \sim \mathcal{U}(-a,\ a)$ & 简单，但未考虑\underline{\hspace{2.0cm}} \\[1.2em]
            Xavier 初始化 & $\sigma = \sqrt{\dfrac{2}{n_{\text{in}}+n_{\text{out}}}}$ & 适合\underline{\hspace{2.5cm}}激活函数 \\[1.5em]
            He 初始化 & $\sigma = \sqrt{\dfrac{2}{n_{\text{in}}}}$ & 专为\underline{\hspace{2.5cm}}激活函数设计 \\
            \bottomrule
        \end{tabular}
        \end{center}

        \item \textbf{为什么常数初始化（全0或全同值）会让网络失效}：从"对称性"角度解释，并说明偏置（bias）为什么通常可以初始化为0。

        \item \textbf{Xavier 初始化的推导思路}：Xavier 初始化的目标是让每一层输出的方差与输入方差相同。在不使用激活函数的假设下，若输入有 $n_{\text{in}}$ 个特征，则将权重初始化为 $\mathcal{N}(0, 1/n_{\text{in}})$ 可以保持方差不变。解释为什么 Xavier 改用 $\sqrt{2/(n_{\text{in}}+n_{\text{out}})}$ 同时考虑了输入和输出维度。

        \item \textbf{He 初始化与 ReLU 的配合}：He 初始化的方差比 Xavier 大了约一倍（分母只有 $n_{\text{in}}$，没有 $n_{\text{out}}$）。解释为什么 ReLU 需要更大的初始方差，以及如果对 ReLU 网络使用 Xavier 初始化会发生什么。
    \end{enumerate}

    \vspace{0.2cm}
    {\color{gray}\small\textit{来源溯源：[文档第 5 章 5.3 参数初始化]}}
\end{problembox}

\begin{beginnerbox}[小白必读：初始化为什么不能"随便设"？]
    \begin{itemize}
        \item \textbf{初始化过大}：权重绝对值很大，激活函数（Sigmoid/Tanh）输入进入饱和区，导数接近0，梯度消失，网络无法更新。
        \item \textbf{初始化过小}：每层输出方差不断缩小，激活值趋近于0，信号在前向传播中逐渐"消亡"，网络同样无法训练。
        \item \textbf{全同初始化的对称性陷阱}：如果同一层所有神经元权重相同，正向传播时输出相同，反向传播时梯度相同，更新后仍然相同——这层退化成"一个神经元"，无论多宽都形同虚设。
        \item \textbf{Xavier/He 的共同目标}：让信号在前向传播中不放大、不缩小地流过每一层，保持激活值的方差在合理范围内。
    \end{itemize}
\end{beginnerbox}

\begin{solutionbox}[参考答案与详细解析]
    \textbf{(1) 六种初始化方法对比表：}

    \begin{center}
    \begin{tabular}{p{2.8cm} p{4.5cm} p{3.8cm}}
        \toprule
        初始化方法 & 规则 & 适用/注意事项 \\
        \midrule
        常数初始化 & 所有权重 $= c$ & 仅限偏置初始化为0；权重禁用 \\[0.8em]
        秩初始化 & 低秩矩阵分解确定初值 & 特定结构网络，不常见 \\[0.8em]
        正态分布初始化 & $\mathcal{N}(0, \sigma^2)$ & $\sigma$ 过大易梯度爆炸，过小易梯度消失 \\[0.8em]
        均匀分布初始化 & $\mathcal{U}(-a, a)$ & 简单，未考虑网络深度/宽度 \\[0.8em]
        Xavier 初始化 & $\sigma=\sqrt{2/(n_{\text{in}}+n_{\text{out}})}$ & 适合 Sigmoid/Tanh 激活函数 \\[0.8em]
        He 初始化 & $\sigma=\sqrt{2/n_{\text{in}}}$ & 专为 ReLU 及其变体设计 \\
        \bottomrule
    \end{tabular}
    \end{center}

    \textbf{(2) 常数初始化失效的原因（对称性问题）：}

    若同一层所有权重 $w_{ij}=c$，则每个神经元接收到完全相同的加权输入，输出完全相同（正向对称）；反向传播时各神经元梯度也完全相同（梯度对称），参数更新方向和幅度一致，更新后权重仍然相同。整层退化为等效的"一个神经元"，无法学习不同特征。

    \textbf{偏置可以初始化为0}：偏置不与同层神经元"对称连接"，不同神经元的偏置独立更新，初始全0不会导致对称性问题，且有助于在训练初期让激活函数在线性区域工作。

    \textbf{(3) Xavier 同时考虑输入和输出维度的原因：}

    在前向传播中，输入维度 $n_{\text{in}}$ 决定了方差的放大倍数；在反向传播中，输出维度 $n_{\text{out}}$ 决定了梯度方差的放大倍数。Xavier 取两者的调和平均 $(n_{\text{in}}+n_{\text{out}})/2$ 作为分母，使得前向和反向传播的方差都尽可能保持不变，是一种兼顾正向和反向信号传播的折中方案。

    \textbf{(4) He 初始化与 ReLU 的配合：}

    ReLU 对所有负输入输出0，相当于在每次前向传播中"关掉"约一半的神经元，使有效传播的信号方差缩小约一半。若使用 Xavier 初始化（设计时假设激活函数不截断信号），每层输出方差会因 ReLU 的截断效应持续减半，导致深层信号趋近于0。He 初始化将方差放大一倍（$2/n_{\text{in}}$），恰好补偿 ReLU 使方差减半的效果，使各层激活值的方差保持稳定。

    \begin{keybox}[初始化方法选择口诀]
        \begin{itemize}
            \item \textbf{Sigmoid / Tanh} $\Rightarrow$ \textbf{Xavier 初始化}（方差 $= 2/(n_\text{in}+n_\text{out})$）
            \item \textbf{ReLU / Leaky ReLU} $\Rightarrow$ \textbf{He 初始化}（方差 $= 2/n_\text{in}$）
            \item \textbf{权重}：永远不要全常数初始化；\textbf{偏置}：通常初始化为0
        \end{itemize}
    \end{keybox}
\end{solutionbox}


%% ============================================================
\Needspace{5\baselineskip}
\begin{problembox}[Q5-05\quad Batch Normalization 代码追踪与行为分析]
    \textbf{题目描述：}\\
    文档第 5.4.1 节介绍了批量标准化（Batch Normalization，BN）。BN 是现代深度网络的"标配"组件，但其训练和推理行为的差异是一个常见的踩坑点。

    \vspace{0.3cm}
    以下是一段含 BN 层的网络代码：
    \begin{lstlisting}[language=Python]
import torch
import torch.nn as nn

class BNNet(nn.Module):
    def __init__(self):
        super().__init__()
        self.fc1    = nn.Linear(4, 8)
        self.bn1    = nn.BatchNorm1d(8)
        self.relu   = nn.ReLU()
        self.fc2    = nn.Linear(8, 2)

    def forward(self, x):
        x = self.fc1(x)      # 步骤1
        x = self.bn1(x)      # 步骤2
        x = self.relu(x)     # 步骤3
        x = self.fc2(x)      # 步骤4
        return x

model = BNNet()
x_train = torch.randn(32, 4)  # batch_size=32, features=4
    \end{lstlisting}

    \textbf{问题：}
    \begin{enumerate}[(1)]
        \item \textbf{BN 的计算过程}：步骤2中，\texttt{bn1} 对输入 $\mathbf{z}$（形状 \texttt{[32, 8]}）进行了哪些操作？写出完整的数学公式，并说明每一步的作用。

        \item \textbf{训练 vs 推理的差异}：\texttt{nn.BatchNorm1d} 在训练阶段和推理阶段的行为有何不同？推理时使用的均值和方差从哪里来？

        \item \textbf{BN 的位置}：BN 层放在线性层\textbf{之后}、激活函数\textbf{之前}是推荐做法，但也有观点认为放在激活函数之后效果更好。分别说明两种方式的理由，以及实践中的推荐选择。

        \item \textbf{BN 的可学习参数}：\texttt{bn1} 中有两个可学习参数 $\gamma$（缩放）和 $\beta$（平移）。解释为什么归一化之后还要进行缩放和平移？若没有这两个参数，BN 会丧失什么能力？

        \item 以下代码中，哪一行会导致 BN 行为异常？原因是什么？
        \begin{lstlisting}[language=Python]
model.train()
for x_batch, y_batch in dataloader:
    pred = model(x_batch)
    loss = criterion(pred, y_batch)
    loss.backward()
    optimizer.step()

# 推理阶段（有问题吗？）
pred_test = model(x_test)   # 行X：未切换推理模式
        \end{lstlisting}
    \end{enumerate}

    \vspace{0.2cm}
    {\color{gray}\small\textit{来源溯源：[文档第 5 章 5.4.1 Batch Normalization 批量标准化]}}
\end{problembox}

\begin{beginnerbox}[小白必读：BN 解决了什么问题？]
    \begin{itemize}
        \item \textbf{内部协变量偏移（Internal Covariate Shift）}：每一层的输入分布会随着上一层参数的更新而变化，导致每一层都要不断适应新的分布，训练极不稳定。BN 在每层输入处强制将激活值标准化为均值0、方差1（再通过 $\gamma$、$\beta$ 调整），使分布保持相对稳定。
        \item \textbf{允许使用更大学习率}：没有 BN 时，过大的学习率容易让激活值进入饱和区，导致梯度消失；BN 将激活值控制在合理范围，即使学习率较大也不容易爆炸。
        \item \textbf{轻微的正则化效果}：BN 利用 mini-batch 的均值和方差来标准化，引入了一定的随机性（不同 batch 的统计量不同），类似于 Dropout 的噪声效果，有轻微防止过拟合的作用。
    \end{itemize}
\end{beginnerbox}

\begin{solutionbox}[参考答案与详细解析]
    \textbf{(1) BN 的完整计算过程：}

    设输入 $\mathbf{z}$ 形状为 \texttt{[32, 8]}（32个样本，8个特征），BN 对每个特征维度（每列）独立进行：

    \textbf{第一步：计算 mini-batch 均值}
    \[
    \mu_j = \frac{1}{N}\sum_{i=1}^{N} z_{ij}, \quad j=1,\ldots,8
    \]
    \textbf{第二步：计算 mini-batch 方差}
    \[
    \sigma_j^2 = \frac{1}{N}\sum_{i=1}^{N}(z_{ij}-\mu_j)^2
    \]
    \textbf{第三步：归一化}
    \[
    \hat{z}_{ij} = \frac{z_{ij}-\mu_j}{\sqrt{\sigma_j^2+\epsilon}}
    \]
    \textbf{第四步：缩放与平移}（$\gamma_j$、$\beta_j$ 为可学习参数）
    \[
    y_{ij} = \gamma_j \hat{z}_{ij} + \beta_j
    \]

    \textbf{(2) 训练 vs 推理的差异：}

    \begin{itemize}
        \item \textbf{训练阶段}：使用当前 mini-batch 的均值 $\mu$ 和方差 $\sigma^2$ 进行归一化，同时用指数移动平均维护全局的 \texttt{running\_mean} 和 \texttt{running\_var}；
        \item \textbf{推理阶段}：不再依赖当前 batch（推理时可能 batch\_size=1），改用训练过程中积累的 \texttt{running\_mean} 和 \texttt{running\_var} 进行归一化，保证输出的确定性。
    \end{itemize}

    \textbf{(3) BN 的位置：}

    \begin{itemize}
        \item \textbf{线性层之后、激活函数之前}（原始论文推荐）：归一化的是线性变换的输出（未激活），数值分布最规则，归一化效果最稳定；
        \item \textbf{激活函数之后}：有研究表明在某些任务（如残差网络）中效果更好，但并非通用结论。
    \end{itemize}

    \textbf{实践推荐}：从"Linear → BN → ReLU"开始，若效果不理想再尝试其他顺序。

    \textbf{(4) $\gamma$ 和 $\beta$ 的必要性：}

    归一化后所有层的激活值都被强制压缩到均值0、方差1的范围，丧失了原本可能很有用的分布形态（比如某层本来激活值应该集中在正值区域）。$\gamma$ 和 $\beta$ 是可学习的，网络可以通过训练自动决定"需要多大范围、需要偏移多少"，从而在标准化的稳定性和表达能力之间取得最优平衡。若去掉这两个参数，网络的表达能力会受到严格限制，极端情况下甚至退化（当 $\gamma=1$，$\beta=0$ 恰好是纯归一化，是 BN 能学到的一种特殊情况）。

    \textbf{(5) 代码问题分析：}

    \textbf{行X 会导致 BN 行为异常}。训练结束后模型仍处于 \texttt{train()} 模式，BN 层在推理时会继续用 \texttt{x\_test} 这一个 batch 的统计量进行归一化，而非使用训练积累的 \texttt{running\_mean}/\texttt{running\_var}。若测试集只有少量样本（甚至单个样本），统计量极不稳定，输出严重失真。

    \textbf{修复方法}：在推理前调用 \texttt{model.eval()}：
    \begin{lstlisting}[language=Python]
model.eval()
with torch.no_grad():
    pred_test = model(x_test)
    \end{lstlisting}
\end{solutionbox}


%% ============================================================
\Needspace{5\baselineskip}
\begin{problembox}[Q5-06\quad 正则化三剑客：权值衰减、Dropout 与 BatchNorm]
    \textbf{题目描述：}\\
    文档第 5.4 节介绍了三种常见正则化手段。过拟合是深度学习实践中最常见的问题，三种方法从不同角度对其发起攻击。

    \vspace{0.3cm}
    \textbf{场景}：你训练了一个5层全连接网络用于图像分类，发现训练集准确率 98\%，验证集准确率仅 72\%，明显过拟合。

    \vspace{0.3cm}
    \textbf{问题：}
    \begin{enumerate}[(1)]
        \item \textbf{权值衰减（L2 正则化）}：在损失函数中加入参数的 L2 范数惩罚项：
        \[
        \mathcal{L}_{\text{total}} = \mathcal{L}_{\text{task}} + \lambda \sum_{l}\|W^{(l)}\|_F^2
        \]
        解释这个惩罚项如何在数学上抑制过拟合，以及 $\lambda$ 过大或过小分别会带来什么问题。

        \item \textbf{Dropout}：训练时以概率 $p$ 随机将神经元输出置零，推理时关闭 Dropout。阅读以下代码，指出其中的错误并改正：
        \begin{lstlisting}[language=Python]
class MyNet(nn.Module):
    def __init__(self):
        super().__init__()
        self.fc1     = nn.Linear(128, 64)
        self.dropout = nn.Dropout(p=0.5)
        self.fc2     = nn.Linear(64, 10)

    def forward(self, x):
        x = torch.relu(self.fc1(x))
        x = self.dropout(x)
        x = self.fc2(x)
        return x

model = MyNet()
# 训练循环
model.train()
for x_batch, y_batch in train_loader:
    pred = model(x_batch)
    loss = criterion(pred, y_batch)
    loss.backward()
    optimizer.step()

# 验证
for x_val, y_val in val_loader:
    pred_val = model(x_val)   # 错误：此处 Dropout 仍在生效
    val_loss = criterion(pred_val, y_val)
        \end{lstlisting}

        \item \textbf{三种方法的作用机制对比}：填写下表，从"如何抑制过拟合"的角度分析三种方法：

        \begin{center}
        \begin{tabular}{p{2.5cm} p{4.5cm} p{4.0cm}}
            \toprule
            方法 & 抑制过拟合的机制 & 额外副作用/优点 \\
            \midrule
            权值衰减 & 惩罚\underline{\hspace{2cm}}权重，迫使模型保持简单 & \underline{\hspace{3.5cm}} \\[1.5em]
            Dropout & 随机关闭神经元，迫使网络不依赖\underline{\hspace{1.5cm}} & 相当于集成\underline{\hspace{1.5cm}}个不同结构的子网络 \\[1.5em]
            BatchNorm & 通过归一化引入\underline{\hspace{1.5cm}}噪声，间接正则化 & 主要功能是\underline{\hspace{2cm}}，正则化是副产品 \\
            \bottomrule
        \end{tabular}
        \end{center}

        \item \textbf{组合使用策略}：针对上述过拟合场景，给出一套包含这三种方法的综合解决方案，说明各参数的推荐取值范围和注意事项。它们可以同时使用吗？是否存在冲突？
    \end{enumerate}

    \vspace{0.2cm}
    {\color{gray}\small\textit{来源溯源：[文档第 5 章 5.4 正则化]}}
\end{problembox}

\begin{beginnerbox}[小白必读：过拟合是什么？为什么深度网络特别容易过拟合？]
    \begin{itemize}
        \item \textbf{过拟合的本质}：模型参数太多、数据太少时，模型有"余力"去记住训练数据中的噪声（如某张猫图有一个特殊的划痕，模型记住"划痕=猫"），而不是学习真正的规律。测试集上没有这些噪声，泛化失败。
        \item \textbf{深度网络为何更容易过拟合}：参数量往往是数据量的数十乃至数百倍。一个1000万参数的网络，理论上可以"完美记忆"数万张图片，但这种记忆对新数据毫无用处。
        \item \textbf{正则化的共同目标}：\textbf{限制模型的"自由度"}，让它不得不学习泛化规律，而不是死记硬背训练集。
    \end{itemize}
\end{beginnerbox}

\begin{metaphorbox}[生活比喻：三种正则化像三种考试备考策略]
    \begin{itemize}
        \item \textbf{权值衰减} = 限制复习时间（参数不能太大），逼你把时间用在最重要的知识点上，而不是背每道题的奇怪细节；
        \item \textbf{Dropout} = 复习时随机遮住一些笔记（随机关闭神经元），逼你从不同角度理解同一知识点，不依赖某一张特定的笔记页；
        \item \textbf{BatchNorm} = 定期对笔记进行整理归档（归一化激活值），让知识体系保持稳定，不让某一个知识点"权重"过大，顺带有轻微防死记的效果。
    \end{itemize}
\end{metaphorbox}

\begin{solutionbox}[参考答案与详细解析]
    \textbf{(1) 权值衰减的数学机制：}

    惩罚项 $\lambda\|W\|^2$ 对大权重施加额外惩罚，促使优化器倾向于选择更小的权重。更新时梯度变为：
    \[
    \frac{\partial \mathcal{L}_{\text{total}}}{\partial W} = \frac{\partial \mathcal{L}_{\text{task}}}{\partial W} + 2\lambda W
    \]
    等效于每次更新前将权重乘以 $(1-2\lambda\eta)$ 进行缩放（"衰减"）。权重越小，模型对输入的敏感性越低，对训练集噪声的过度拟合能力越弱。

    \begin{itemize}
        \item $\lambda$ \textbf{过大}：权重被压缩过度，模型欠拟合，连训练集规律都学不到；
        \item $\lambda$ \textbf{过小}：正则化效果微弱，过拟合依然存在。
    \end{itemize}
    常用范围：$\lambda \in [10^{-4}, 10^{-2}]$，具体由验证集表现决定。

    \textbf{(2) Dropout 代码错误与修正：}

    \textbf{错误}：验证循环中未切换到 \texttt{eval()} 模式，Dropout 仍以 $p=0.5$ 随机关闭神经元，导致验证集的 Loss 和准确率每次都不一样，无法正确评估模型性能。

    \textbf{修正}：
    \begin{lstlisting}[language=Python]
# 验证前切换模式，并关闭梯度计算
model.eval()
with torch.no_grad():
    for x_val, y_val in val_loader:
        pred_val = model(x_val)   # Dropout 已关闭
        val_loss = criterion(pred_val, y_val)
# 验证结束后切回训练模式
model.train()
    \end{lstlisting}

    \textbf{(3) 三种方法的作用机制对比表：}

    \begin{center}
    \begin{tabular}{p{2.5cm} p{4.5cm} p{3.8cm}}
        \toprule
        方法 & 抑制过拟合的机制 & 额外副作用/优点 \\
        \midrule
        权值衰减 & 惩罚过大权重，迫使模型保持简单 & 超参数 $\lambda$ 需调节，可能欠拟合 \\[0.8em]
        Dropout & 随机关闭神经元，迫使网络不依赖某些特定特征 & 相当于集成指数级个不同结构的子网络 \\[0.8em]
        BatchNorm & 通过归一化引入 mini-batch 噪声，间接正则化 & 主要功能是加速训练/稳定梯度，正则化是副产品 \\
        \bottomrule
    \end{tabular}
    \end{center}

    \textbf{(4) 综合解决方案：}

    三种方法\textbf{可以同时使用}，通常不存在根本冲突，但需注意：
    \begin{itemize}
        \item \textbf{BatchNorm + Dropout 的顺序}：通常先 BN 后 Dropout，且若网络已有 BN，Dropout 的 $p$ 值可适当减小（如0.3），因为 BN 本身已提供一定正则化；
        \item \textbf{推荐配置}（针对本题场景）：
        \begin{lstlisting}[language=Python]
# 每个隐藏层结构：Linear -> BN -> ReLU -> Dropout
nn.Linear(in, out),
nn.BatchNorm1d(out),
nn.ReLU(),
nn.Dropout(p=0.3),       # 适度 Dropout，不要 p=0.5 过强
        \end{lstlisting}
        \item \textbf{权值衰减}：在优化器中加入 \texttt{weight\_decay=1e-4}；
        \item \textbf{注意事项}：三种方法都会使训练 Loss 高于"没有正则化"时的值，这是正常的——目标是验证集 Loss，不是训练集 Loss。正则化强度过大会导致欠拟合，需通过验证集调参。
    \end{itemize}

    \begin{keybox}[过拟合处置优先级建议]
        \begin{itemize}
            \item \textbf{第一步}：增加数据量（最根本的解决方案）；
            \item \textbf{第二步}：加入 Dropout（$p=0.3\sim0.5$）；
            \item \textbf{第三步}：加入权值衰减（$\lambda=10^{-4}\sim10^{-3}$）；
            \item \textbf{第四步}：加入 BatchNorm（同时加速训练）；
            \item \textbf{第五步}：考虑减小模型规模（减少层数/宽度）。
        \end{itemize}
    \end{keybox}
\end{solutionbox}

\begin{appbox}[现实应用场景]
    \textbf{这些"技巧"在工业界的实际地位：}\\
    本章所有内容在工业级深度学习中都是标配，而非可选项。以 GPT 系列大模型为例：它使用 Adam 优化器（带权值衰减的 AdamW 变体）、余弦退火学习率调度、He 初始化、LayerNorm（BatchNorm 的序列友好变体）。ResNet 之所以能训练到上百层，核心就是 BatchNorm 解决了深层网络的梯度问题，以及残差连接提供了额外的梯度通路。每一个本章的"技巧"，都是研究者在反复踩坑之后总结出来的，理解它们背后的原理，才能在遇到新问题时知道该用哪把钥匙。
\end{appbox}

% ===== 本章小结 =====
\section{本章小结：把"技巧"变成本能}

学完这 6 道题，如果还觉得"内容繁杂、难以归纳"，试试把整章内容浓缩成这几句画面：

\begin{itemize}
    \item \textbf{梯度消失/爆炸} = 深层网络的"信号失真"问题；换 ReLU、用合适初始化、加 BatchNorm，三管齐下。
    \item \textbf{优化器进化链}：SGD $\to$ Momentum（加惯性）$\to$ AdaGrad（自适应步长）$\to$ RMSProp（修复学习率归零）$\to$ Adam（全都要）；学习率衰减是所有优化器的好搭档。
    \item \textbf{初始化的铁律}：权重不能全同；Sigmoid/Tanh 用 Xavier，ReLU 用 He；偏置初始化为0。
    \item \textbf{BN 的核心}：训练用 batch 统计量，推理用 running 统计量，推理前必须 \texttt{model.eval()}！
    \item \textbf{正则化三剑客}：权值衰减限制权重大小，Dropout 阻断对特定神经元的依赖，BN 通过归一化噪声正则化——三者可以叠加使用。
    \item \textbf{过拟合的终极诊断}：训练集 Loss 低、验证集 Loss 高 $\Rightarrow$ 过拟合，加正则化；两者都高 $\Rightarrow$ 欠拟合，减正则化或增大模型。
\end{itemize}

\begin{metaphorbox}[给零基础读者的一句总比喻]
    \textbf{如果把训练深度网络比作培养一名运动员：}
    \begin{itemize}
        \item \textbf{梯度消失/爆炸} = 训练信号失真，教练（梯度）的指令传不到运动员耳里，或者喊得太大声把运动员吓到——换一套传令系统（ReLU + 合理初始化）；
        \item \textbf{优化器} = 训练计划：只跟着当天状态练（SGD），还是结合历史成绩制定长期方案（Adam）？
        \item \textbf{初始化} = 运动员的起跑姿势：一开始站对了，才能跑得快；
        \item \textbf{BatchNorm} = 每次训练前的热身和拉伸，让身体保持在最佳状态再开始发力；
        \item \textbf{正则化} = 防止过度专项训练：光练100米短跑（过拟合训练集），遇到障碍跑就废了；要练综合素质（泛化能力），才能在任何比赛中都有好成绩。
    \end{itemize}
\end{metaphorbox}

\end{document}
